\newpage
\chapter{Computability}


\subsection{Church Turing Thesis}

The intuitive computable functions are exactly all functions which can be computed 
by a \emph{Turing Machine}. In other words every computational model can be expressed as such.

\subsection{Textual Substitution}

For a given condition of the form \(\{E[x:=R]\}\) we perform the substitution in the following way 

\[
    (2x + y)[x:= 4y] = 8y + y = 9y
\]

If there are multiple replacements of the form \([\dots], \dots, [\dots]\) we perform those from left to right.

\subsection{Leibniz Inference Rule}

To clarify the notation we will be using from now looks like fraction but it means that assuming 
the statement above is true then the following (bellow) is also true. 

\[
    \frac{X = Y}{E[z:=X] = E[z:= Y]} 
\]

\subsection{Correctness of Algorithms}

The \emph{correctness} of algorithm is determined by the sum of the following properties 
the \emph{partial correctness}, this means the algorithm at each step does the right thing, and 
\emph{termination}, which like the name means that the algorithm terminates. 

\subsubsection{Partial Correctness}

If before running \(P\) the condition \(\varphi\) is true, then after the termination \(P\) should 
\(\psi\) be true; these are the pre- and post-conditions.

\textbf{Example:}

\[
    <x=0> x:= x+1 <x>0>
\]

\subsubsection{Hoare Calculus}

To read the following expression which look like fraction but are not, use the following:

\begin{itemize}
    
    \item \emph{Above}: if this is true 
    
    \item \emph{Below}: then this is also true

\end{itemize}

\emph{Assignation Rule}

If \textbf{x} is a variable, \textbf{t} is an expression without side effects \(\varphi \text{x/t}\) is 
\(\psi\) with all \textbf{x} substituted for \textbf{t}.

\[
    <\varphi [\text{x/t}]> \text{ x} = \text{t; } <\psi>
\]

Note that this rule only has the part below, this means that this rule has no precondition.

\textbf{Example I:} 

\(<\text{y } = 5>\) x = 5; \(<\text{x } = \text{ y}>\)

\textbf{Example II:} 

\(<5 > 4>\) x = 5; \(<\text{x } > 4>\)


\emph{Consequence Rule (Strong pre-condition)} 

Here \(\alpha\) is a stronger condition than \(\phi\). This way our condition \(\phi\) has an extra layer of protection.
\[
    \frac{<\varphi> P <\psi> \quad \alpha \implies \varphi}{ <\alpha> P <\psi>}
\]

\textbf{Example:} 

\[
    \frac{<5 = 5> \text{x } = 5; <\text{x } = 5> \text{ true } \implies 5 = 5}{<\text{true} > x = 5; <\text{x } = 5>}
\]

In code 

\begin{verbatim}
    <true> 
    <5=5> 
    x=5;
    <x=5>
\end{verbatim}

\emph{Consequence Rule (Weak post-condition)} 

This let us make the post-condition weaker.

\[
    \frac{<\varphi> P <\psi> \quad \psi \implies \beta}{ <\varphi> P <\beta>}
\]

\textbf{Example:} 

\[
    \frac{<5 = 5> \text{x } = 5; <\text{x } = 5> \text{ x } = 5 \implies x \ge 5}{<\text{true}> x = 5; <\text{x } \ge 5>}
\]

Note that two guards can only be consecutive if from the one above the one below is implied 

\emph{Sequence Rule}

This rule allows us to compress our conditions.

\[
    \frac{<\psi> P <\psi> \quad <\psi> Q <\beta>}{ <\psi> P Q <\beta>}
\]

\textbf{Example:}

\(<\text{true}>\) 
x = 5;
res = x * x + 6;
\(<\text{res } = 31>\)

\emph{Condition Rule I}

This rule enforces that no matter in which branch of the program  is taken the post-condition  
is true.

\[
    \frac{<\phi \land B > P <\psi> \quad \phi \land \neg B \implies \psi}{ <\phi> \text{ if } (B) \{P\} <\psi>}
\]

\textbf{Example:}

\begin{verbatim}
<true>
<y=y>
res=y
<res=y>

if (x > y) {
    <res = y \land x > y>
    <x = max(x,y)>
    res = x 
    <res = max(x,y)>   
}
<res = max(x,y)>
\end{verbatim}

\emph{Condition Rule II}

Similar to the previous case, we want to make sure that independent of the branch 
we have a correct post-condition.

\[
    \frac{<\phi \land B > P <\psi> \quad <\phi \land \neg B> Q <\psi>}{ <\phi> \text{ if } (B) \text{ else } \{Q\} <\psi>}
\]

\textbf{Example:}

\begin{verbatim}
<true>
<x=x>
res=x
<res=x>

if (x < 0) {
    <true \land x < 0>
    <-x = |x|>
    res = -x 
    <res = |x|>   
} else {
    <true \land -x < 0 >
    <x = |x|>
    res = x 
    <res = |x>
}
<res = |x|>
\end{verbatim}

\emph{Loop Rule}

In this rule we look for some \emph{invariant} \(\phi\) which before and after the loop is true.

\[
    \frac{<\phi \land B > P <\phi>}{ <\phi> \text{ while } (B) \{Q\} <\phi \land \neg B>}
\]

Note that from the guard before the loop \(<\alpha>\), we have \(\alpha \implies \phi\)
Finding the invariant can be tricky, but it commonly involves finding an integer expression, also some formula which 
holds for the entire loop.

\textbf{Example:}

In this case the invariant is \(i! res = n!\)

\begin{verbatim}
    <true>
    i = n;
    res = 1;
    <i = n \land res = 1>
    <i! res = n!>
    while (i > 1) {
    <i! res = n! \land i > 1>
    <(i - 1)! res i = n!>
        res = res * i;
    <(i - 1)! res = n!>
        i = i - 1;
    <i! res = n!>
    }
    <\phi \land \neg i > 1>
    <res = n!>
\end{verbatim}

\subsection{Termination}

For each loop \textbf{while} \((B) \{P\}\) we have to find some \textbf{int} expression \(V\) 
called the \emph{variant} such that 

\[
    B \implies V \ge 0 \text{ and } <V = m \land B> P <V < m> >
\]

In simpler words, we look for an integer expression which continuously decreases.

\textbf{Example:}

\begin{verbatim} 
    <i = m \land i > 1>
    while ( i > q) {

    <i - 1 < m >
        res  = res * i;
    <i -1 < m>
        i = i - 1;;
    <i < m>
    }
\end{verbatim}


